\chapter{Introduction}
\label{sec:introduction}

\section{Background and Motivation}
Information access systems help users locate material that is relevant to a
question or task. A search interface may return a ranked list of documents,
while another interface may combine retrieved passages into a short answer.
In either case, evaluation begins with a clear account of the user need and
the evidence that would satisfy it. A system that retrieves useful material
for one type of question may perform differently when the vocabulary, topic,
or expected level of detail changes.

This illustrative thesis considers how to compare two retrieval configurations
under a shared evaluation procedure. The purpose of the example is to provide
realistic paragraph lengths and document structure for checking the template.
It does not report a completed experiment. Any statements about the example
systems describe the proposed study design, and the numerical values shown
later are invented. Keeping this distinction explicit also demonstrates how
a thesis can separate an experimental plan from observed results.

\section{Research Questions}
The example asks whether an additional reranking stage improves the usefulness
of the first few retrieved documents, and whether any improvement is consistent
across different query groups. These are separate questions. An increase in an
overall mean does not establish that every group benefits, because the mean
can combine gains on some queries with losses on others. The evaluation must
therefore retain query-level observations as well as summary values.

The main comparison uses the same collection, queries and relevance judgements
for both configurations. Only the ranking procedure changes. This restriction
makes the comparison easier to interpret, although it does not eliminate all
sources of uncertainty. Later chapters discuss judgement coverage, parameter
selection and the limits of a small evaluation collection.

% Demonstration-only break: expose an ordinary continuation page.
\newpage
\section{Scope and Contributions}
\label{sec:scope}
The intended contribution of this example study is an explicit and reproducible
comparison procedure. A useful report should explain how the input collection
was prepared, how queries were selected, which settings were held constant,
and how per-query scores were aggregated. The reader should be able to trace
a value in a results table back to the observations and choices that produced
it. A single headline score is not a substitute for this account.

Three outputs are envisaged. The first is a documented baseline configuration.
The second is a controlled variant that introduces a reranking stage. The third
is an analysis that considers both average performance and variation across
queries. These outputs are described as a plan, not as completed research
contributions. In a real thesis, this section should be rewritten to reflect
what was actually implemented and established by the evidence.

\subsection{Assumptions}
The example assumes that each query has a fixed information need and that the
available judgements are used consistently for both systems. It also assumes
that the collection remains unchanged during the experiment. These assumptions
simplify the comparison, but they limit the conclusions: an offline experiment
does not directly establish how people would use the system in a live setting.
The distinction should remain visible when reporting the results.

\subsection{Reading the Example}
The document includes both numbered chapters and unnumbered back matter.
Tables and figures have independent numbering, and references to them are
resolved automatically. For example, Chapter~\ref{sec:experiments} describes the
experimental design, while Appendix~\ref{app:details} lists the information
needed to reproduce it. These references should remain correct after material
is added or removed.

% A second continuation page exposes the other odd/even running header.
\newpage
\section{Organisation of the Thesis}
Chapter~\ref{sec:related_work} explains how to organise related work around the
decisions made in an evaluation. Chapter~\ref{sec:approach} then defines the
illustrative comparison and its aggregation rule. The experimental chapter
provides example tables, and Chapter~\ref{sec:results_and_discussions} shows how
to discuss an invented result without presenting it as an empirical finding.
The final chapters distinguish conclusions from possible extensions.

This organisation places the study design before the results. That order helps
the reader determine which decisions were part of the original procedure and
which interpretations arose afterwards. A complete thesis may need additional
chapters for implementation, datasets or user studies. The template does not
require every project to follow the particular chapter structure used here.

\section{Layout Inspection}
On this continuation page, the running header should identify the current
chapter or section, according to page parity. The footer should display the
page number in the same horizontal position as on the chapter-opening page.
The opening page itself should have no running header and no header rule.
This first chapter has three pages so that both ordinary header positions can
be inspected without changing the source.

The next chapter begins on the immediately following page. Its page number
may be even or odd; no empty page is required between chapters. The forced
page breaks within these sample chapters are solely for demonstrating page
styles. During normal writing, paragraphs flow automatically and each
\verb|\chapter| command starts a new page on its own.
